% ---------------------------------------------------------------------------
% SdS - TP2: Autómatas celulares y bandadas de agentes autopropulsados
% Presentación oral: 13 minutos.
%
% Formato según material/catedra/practica/GuiaPresentaciones.pdf y las
% correcciones devueltas a otros grupos, resumidas en
% restricciones-por-seccion.md.
%
% Los recuadros grises son ranuras de contenido pendiente. Cada uno indica el
% archivo esperado en figuras/ y los datos que debe traer. Reemplazar
% \slot{...} por \includegraphics[width=\textwidth]{figuras/<archivo>}.
% Lo marcado con \pend{...} son valores todavía no medidos.
% ---------------------------------------------------------------------------
\documentclass[aspectratio=169,10pt]{beamer}

\usetheme{Warsaw}
\useoutertheme{miniframes}
\setbeamertemplate{navigation symbols}{}
% miniframes trae el pie vacío: hay que pedir el número de diapositiva explícitamente.
\setbeamertemplate{footline}[frame number]

\usepackage[utf8]{inputenc}
\usepackage[T1]{fontenc}
\usepackage[spanish,es-noquoting]{babel}
\usepackage{amsmath,amssymb,bm}
\usepackage{siunitx}
\usepackage{graphicx}
\usepackage{booktabs}
\usepackage{xcolor}
\usepackage{url}

\sisetup{exponent-product=\times}

% Ranura de contenido pendiente: #1 ancho, #2 alto, #3 descripción
\newcommand{\slot}[3]{%
  \fbox{\parbox[c][#2][c]{#1}{\centering\footnotesize\color{gray}#3}}}
% Valor pendiente de medición
\newcommand{\pend}[1]{{\color{red}\textbf{[#1]}}}
% Bloque de parámetros fijos, al costado de la figura y fuera de ella
\newcommand{\params}[1]{{\scriptsize #1}}
% Link a animación: debe verse la URL completa
\newcommand{\anim}[1]{{\small\texttt{#1}}}

\title[SdS TP2]{Bandadas de agentes autopropulsados}
\subtitle{Modelo estándar y modelo de votante}
\author{\pend{Integrante 1} \and \pend{Integrante 2} \and \pend{Integrante 3}}
\institute{Simulación de Sistemas — Grupo \pend{XX}, Comisión \pend{S}}
\date{4 de septiembre de 2026}

\begin{document}

\begin{frame}[plain]
  \titlepage
\end{frame}

% 
\section*{Introducción}
\begin{frame}\sectionpage\end{frame}

\begin{frame}{Sistema real}
  \begin{columns}[T]
    \begin{column}{0.52\textwidth}
      \slot{\textwidth}{4.4cm}{figuras/01-sistema-real.png\\[2pt] foto o esquema de una bandada real}
    \end{column}
    \begin{column}{0.44\textwidth}
      \begin{itemize}
        \item Agentes autopropulsados a rapidez constante.
        \item Sin líder ni control central: solo interacción local.
        \item Orden colectivo emergente frente a perturbaciones.
      \end{itemize}
    \end{column}
  \end{columns}
\end{frame}

\begin{frame}{Modelo matemático}
  Cada agente $i$ tiene posición $\bm{x}_i$ y orientación $\theta_i$, con $|\bm{v}_i| = v$ constante:
  \begin{align}
    \bm{x}_i(t+\Delta t) &= \bm{x}_i(t) + \bm{v}_i(t)\,\Delta t, \\
    \theta_i(t+\Delta t) &= \Theta_i(t) + \Delta\theta,
  \end{align}
  donde $\Delta\theta$ es uniforme en $[-\eta/2,\ \eta/2]$ y $\Theta_i$ es la dirección que la
  regla de interacción asigna a partir de las vecinas dentro del radio $r_c$.
\end{frame}

\begin{frame}{Reglas de interacción}
  \begin{columns}[T]
    \begin{column}{0.48\textwidth}
      \textbf{Estándar}\\[4pt]
      Promedia las direcciones de todas sus vecinas, incluida ella misma:
      \begin{equation}
        \Theta_i = \operatorname{atan2}\!\big(\langle \sin\theta \rangle_{r_c},\ \langle \cos\theta \rangle_{r_c}\big)
      \end{equation}
      {\scriptsize Vicsek \emph{et al.}, \emph{Phys.\ Rev.\ Lett.}, 1995}
    \end{column}
    \begin{column}{0.48\textwidth}
      \textbf{Votante}\\[4pt]
      Elige una sola vecina al azar y copia su dirección:
      \begin{equation}
        \Theta_i = \theta_j, \qquad j \sim \mathcal{U}(\mathcal{V}_i)
      \end{equation}
      {\scriptsize Loscar, Baglietto y Vazquez, \emph{Phys.\ Rev.\ E}, 2021}
    \end{column}
  \end{columns}
  \vspace{0.6em}
  \centering\slot{0.5\textwidth}{2.0cm}{figuras/02-esquema-reglas.pdf\\[2pt] misma vecindad, las dos reglas lado a lado}
\end{frame}

% ===========================================================================
\section*{Implementación}
\begin{frame}\sectionpage\end{frame}

\begin{frame}{Arquitectura del motor de simulación}
  \begin{columns}[T]
    \begin{column}{0.56\textwidth}
      \slot{\textwidth}{4.4cm}{figuras/03-arquitectura.pdf\\[2pt] diagrama de módulos o UML del motor\\ (sin post-proceso ni formato de archivos)}
    \end{column}
    \begin{column}{0.40\textwidth}
      \begin{itemize}\small
        \item \pend{Módulo de estado / partículas}
        \item \pend{Grilla CIM con contorno periódico}
        \item \pend{Regla de actualización (estándar / votante)}
        \item \pend{Integrador temporal}
      \end{itemize}
    \end{column}
  \end{columns}
\end{frame}

\begin{frame}{Paso de simulación}
  \begin{columns}[T]
    \begin{column}{0.56\textwidth}
      \slot{\textwidth}{4.4cm}{figuras/04-pseudocodigo.pdf\\[2pt] pseudocódigo del paso $t \to t+\Delta t$:\\ construir grilla, buscar vecinas,\\ actualizar $\theta$, actualizar $\bm{x}$, aplicar PBC}
    \end{column}
    \begin{column}{0.40\textwidth}
      \begin{itemize}\small
        \item Actualización sincrónica: todas las $\theta$ nuevas se calculan sobre el estado en $t$.
        \item Vecinas por CIM con $M$ celdas por lado, $L/M \geq r_c$.
        \item Distancias con imagen mínima.
      \end{itemize}
    \end{column}
  \end{columns}
\end{frame}

% ===========================================================================
\section*{Simulaciones}
\begin{frame}\sectionpage\end{frame}

\begin{frame}{Sistema particular}
  \begin{columns}[T]
    \begin{column}{0.46\textwidth}
      \slot{\textwidth}{4.6cm}{figuras/05-esquema-caja.pdf\\[2pt] caja con contorno periódico.\\ $L$ y $r_c$ acotados \emph{sobre el dibujo},\\ una partícula con su vector velocidad}
    \end{column}
    \begin{column}{0.50\textwidth}
      \params{
      \begin{tabular}{@{}ll@{}}
        \multicolumn{2}{@{}l}{\textbf{Fijos}} \\
        \midrule
        Lado de la caja & $L = 10$ \\
        Contorno & periódico \\
        Radio de interacción & $r_c = 1$ \\
        Rapidez & $v = \num{3e-2}$ \\
        Paso temporal & $\Delta t = 1$ \\
        Pasos por corrida & \pend{t} \\
        \addlinespace
        \multicolumn{2}{@{}l}{\textbf{Variables}} \\
        \midrule
        Densidad & $\rho = 2,\ 4,\ 8$ \\
        Partículas & $N = \rho L^2 = 200,\ 400,\ 800$ \\
        Ruido & $\eta \in [\pend{$\eta_{\min}$},\ \pend{$\eta_{\max}$}]$ \\
        Regla & estándar, votante \\
        \addlinespace
        Realizaciones & $n_r \in \pend{[a, b]}$ \\
        \bottomrule
      \end{tabular}}
    \end{column}
  \end{columns}
\end{frame}

\begin{frame}{Observables}
  \textbf{Polarización}, instante a instante:
  \begin{equation}
    v_a(t) = \frac{1}{N v}\left| \sum_{i=1}^{N} \bm{v}_i(t) \right| \in [0,1]
  \end{equation}

  \textbf{Fracción de la componente gigante}. Un \emph{cluster} es un conjunto de partículas donde
  todo par está unido por una cadena de vecinas a distancia $\le r_c$:
  \begin{equation}
    S(t) = \frac{1}{N}\max_{c\,\in\,\mathcal{C}(t)} |c| \in (0,1]
  \end{equation}

  \vspace{0.4em}
  \textbf{Observables escalares}: $\langle v_a\rangle$ y $\langle S\rangle$ promedian los valores de
  $v_a(t)$ y $S(t)$ sobre los pasos posteriores a $t_0$, el instante en que el sistema alcanza el
  estacionario, y sobre las $n_r$ realizaciones. La barra de error es la dispersión entre realizaciones.
\end{frame}

% ===========================================================================
\section*{Resultados}
\begin{frame}\sectionpage\end{frame}

% --- Estudio 1: ruido, modelo estándar ------------------------------------
\begin{frame}{Respuesta al ruido: modelo estándar}
  \framesubtitle{Dinámica para ruido bajo y ruido alto}
  \begin{columns}[T]
    \begin{column}{0.46\textwidth}
      \centering
      \slot{\textwidth}{3.0cm}{figuras/06-anim-estandar-eta-bajo.png}\\[3pt]
      \anim{\pend{https://youtu.be/xxxxxxxxxxx}}
    \end{column}
    \begin{column}{0.46\textwidth}
      \centering
      \slot{\textwidth}{3.0cm}{figuras/07-anim-estandar-eta-alto.png}\\[3pt]
      \anim{\pend{https://youtu.be/xxxxxxxxxxx}}
    \end{column}
  \end{columns}
  \vspace{0.5em}
  \params{$L=10$ \quad $\rho = \pend{$\rho$}$ \quad $N = \pend{N}$ \quad $v=\num{3e-2}$ \quad $r_c=1$
  \quad $\eta = \pend{$\eta_{\min}$}$ (izq.) y $\pend{$\eta_{\max}$}$ (der.)}
\end{frame}

\begin{frame}{Respuesta al ruido: modelo estándar}
  \framesubtitle{Evolución temporal de la polarización}
  \begin{columns}[T]
    \begin{column}{0.64\textwidth}
      \slot{\textwidth}{4.8cm}{figuras/10-va-vs-t-estandar.pdf\\[2pt] polarización vs.\ tiempo\\ leyenda: solo los valores de $\eta$\\ línea vertical en $t_0$}
    \end{column}
    \begin{column}{0.32\textwidth}
      \params{
      $L = 10$\\
      $\rho = \pend{$\rho$}$, $N = \pend{N}$\\
      $v = \num{3e-2}$, $r_c = 1$\\
      $n_r = \pend{n}$\\[6pt]
      $t_0 = \pend{t}$: a partir de ahí se promedia para obtener $\langle v_a\rangle$.}
    \end{column}
  \end{columns}
\end{frame}

\begin{frame}{Respuesta al ruido: modelo estándar}
  \framesubtitle{Polarización en función del ruido}
  \begin{columns}[T]
    \begin{column}{0.64\textwidth}
      \slot{\textwidth}{4.8cm}{figuras/11-va-vs-eta-estandar.pdf\\[2pt] polarización vs.\ ruido, con barras de error\\ leyenda: solo los valores de $\rho$}
    \end{column}
    \begin{column}{0.32\textwidth}
      \params{
      $L = 10$\\
      $v = \num{3e-2}$, $r_c = 1$\\
      $n_r = \pend{n}$ por punto\\
      Promedio para $t > t_0$}
    \end{column}
  \end{columns}
\end{frame}

\begin{frame}{Respuesta al ruido: modelo estándar}
  \framesubtitle{Evolución temporal de la componente gigante}
  \begin{columns}[T]
    \begin{column}{0.64\textwidth}
      \slot{\textwidth}{4.8cm}{figuras/12-S-vs-t-estandar.pdf\\[2pt] fracción de la componente gigante vs.\ tiempo\\ leyenda: solo los valores de $\rho$\\ línea vertical en $t_0$}
    \end{column}
    \begin{column}{0.32\textwidth}
      \params{
      $L = 10$\\
      $v = \num{3e-2}$, $r_c = 1$\\
      $\eta = \pend{$\eta$}$\\
      $n_r = \pend{n}$}
    \end{column}
  \end{columns}
\end{frame}

\begin{frame}{Respuesta al ruido: modelo estándar}
  \framesubtitle{Componente gigante en función del ruido}
  \begin{columns}[T]
    \begin{column}{0.64\textwidth}
      \slot{\textwidth}{4.8cm}{figuras/13-S-vs-eta-estandar.pdf\\[2pt] componente gigante vs.\ ruido, con barras de error\\ leyenda: solo los valores de $\rho$}
    \end{column}
    \begin{column}{0.32\textwidth}
      \params{
      $L = 10$\\
      $v = \num{3e-2}$, $r_c = 1$\\
      $n_r = \pend{n}$ por punto\\
      Promedio para $t > t_0$}
    \end{column}
  \end{columns}
\end{frame}

% --- Estudio 2: ruido, modelo de votante ----------------------------------
\begin{frame}{Respuesta al ruido: modelo de votante}
  \framesubtitle{Dinámica para ruido bajo y ruido alto}
  \begin{columns}[T]
    \begin{column}{0.46\textwidth}
      \centering
      \slot{\textwidth}{3.0cm}{figuras/08-anim-votante-eta-bajo.png}\\[3pt]
      \anim{\pend{https://youtu.be/xxxxxxxxxxx}}
    \end{column}
    \begin{column}{0.46\textwidth}
      \centering
      \slot{\textwidth}{3.0cm}{figuras/09-anim-votante-eta-alto.png}\\[3pt]
      \anim{\pend{https://youtu.be/xxxxxxxxxxx}}
    \end{column}
  \end{columns}
  \vspace{0.5em}
  \params{$L=10$ \quad $\rho = \pend{$\rho$}$ \quad $N = \pend{N}$ \quad $v=\num{3e-2}$ \quad $r_c=1$
  \quad $\eta = \pend{$\eta_{\min}$}$ (izq.) y $\pend{$\eta_{\max}$}$ (der.)}
\end{frame}

\begin{frame}{Respuesta al ruido: modelo de votante}
  \framesubtitle{Evoluciones temporales}
  \begin{columns}[T]
    \begin{column}{0.64\textwidth}
      \slot{\textwidth}{4.8cm}{figuras/16-evoluciones-votante.pdf\\[2pt] dos paneles: polarización vs.\ tiempo\\ y componente gigante vs.\ tiempo\\ línea vertical en $t_0$}
    \end{column}
    \begin{column}{0.32\textwidth}
      \params{
      $L = 10$\\
      $\rho = \pend{$\rho$}$, $N = \pend{N}$\\
      $v = \num{3e-2}$, $r_c = 1$\\
      $\eta = \pend{$\eta_{\min}$}$ y $\pend{$\eta_{\max}$}$\\
      $n_r = \pend{n}$\\[6pt]
      Mismo criterio de $t_0$ que en el modelo estándar.}
    \end{column}
  \end{columns}
\end{frame}

\begin{frame}{Respuesta al ruido: modelo de votante}
  \framesubtitle{Polarización en función del ruido, comparada con el modelo estándar}
  \begin{columns}[T]
    \begin{column}{0.64\textwidth}
      \slot{\textwidth}{4.8cm}{figuras/17-va-vs-eta-comparacion.pdf\\[2pt] polarización vs.\ ruido, con barras de error\\ leyenda: valores de $\rho$ y regla de interacción}
    \end{column}
    \begin{column}{0.32\textwidth}
      \params{
      $L = 10$\\
      $v = \num{3e-2}$, $r_c = 1$\\
      $n_r = \pend{n}$ por punto\\
      Promedio para $t > t_0$}
    \end{column}
  \end{columns}
\end{frame}

\begin{frame}{Respuesta al ruido: modelo de votante}
  \framesubtitle{Componente gigante en función del ruido, comparada con el modelo estándar}
  \begin{columns}[T]
    \begin{column}{0.64\textwidth}
      \slot{\textwidth}{4.8cm}{figuras/18-S-vs-eta-comparacion.pdf\\[2pt] componente gigante vs.\ ruido, con barras de error\\ leyenda: valores de $\rho$ y regla de interacción}
    \end{column}
    \begin{column}{0.32\textwidth}
      \params{
      $L = 10$\\
      $v = \num{3e-2}$, $r_c = 1$\\
      $n_r = \pend{n}$ por punto\\
      Promedio para $t > t_0$}
    \end{column}
  \end{columns}
\end{frame}

% --- Estudio 3: relación entre observables --------------------------------
\begin{frame}{Relación entre polarización y componente gigante}
  \begin{columns}[T]
    \begin{column}{0.64\textwidth}
      \slot{\textwidth}{4.8cm}{figuras/14-va-vs-S.pdf\\[2pt] polarización vs.\ componente gigante\\ leyenda: valores de $\rho$ y regla de interacción}
    \end{column}
    \begin{column}{0.32\textwidth}
      \params{
      $L = 10$\\
      $v = \num{3e-2}$, $r_c = 1$\\
      Cada punto corresponde a un valor de $\eta$\\
      $n_r = \pend{n}$}
    \end{column}
  \end{columns}
\end{frame}

% --- Estudio 4: costo computacional ---------------------------------------
\begin{frame}{Tiempos de ejecución del CIM}
  \begin{columns}[T]
    \begin{column}{0.60\textwidth}
      \slot{\textwidth}{4.4cm}{figuras/15-tiempos-cim.pdf\\[2pt] tiempo de ejecución (s) vs.\ número de partículas\\ leyenda: TP1 y TP2}
    \end{column}
    \begin{column}{0.36\textwidth}
      \params{
      \begin{tabular}{@{}lcc@{}}
        \toprule
        $N$ & TP1 (s) & TP2 (s) \\
        \midrule
        \pend{N} & \pend{t} & \pend{t} \\
        \pend{N} & \pend{t} & \pend{t} \\
        \pend{N} & \pend{t} & \pend{t} \\
        \bottomrule
      \end{tabular}\\[6pt]
      $L = 10$, $r_c = 1$, $M = \pend{M}$\\
      \pend{n} repeticiones por punto}
    \end{column}
  \end{columns}
\end{frame}

% ===========================================================================
\section*{Conclusiones}
\begin{frame}\sectionpage\end{frame}

\begin{frame}{Conclusiones}
  \begin{itemize}
    \item \pend{Cómo decrece $\langle v_a\rangle$ al aumentar $\eta$ en el modelo estándar, con el rango numérico observado}
    \item \pend{Cómo cambia ese comportamiento con $\rho$ en el modelo estándar}
    \item \pend{En qué se diferencia el modelo de votante del estándar, nombrando el modelo en cada afirmación}
    \item \pend{Qué relación cuantitativa muestran $\langle v_a\rangle$ y $\langle S\rangle$}
    \item \pend{Cómo se comparan los tiempos del CIM con los del TP1}
  \end{itemize}
  \vspace{0.5em}
  {\scriptsize Cada afirmación nombra el modelo al que corresponde y se sostiene en una figura mostrada.}
\end{frame}

\begin{frame}
  \centering\vfill{\Large Muchas gracias}\vfill
\end{frame}

\end{document}
